\documentclass{amsart}
\usepackage[UTF8,fontset=windows]{ctex}
\usepackage[a4paper,margin=24mm]{geometry}
% For dealing with references we use the comment environment
\usepackage{verbatim}
\newenvironment{reference}{\comment}{\endcomment}
%\newenvironment{reference}{}{}
\newenvironment{slogan}{\comment}{\endcomment}
\newenvironment{history}{\comment}{\endcomment}

% For commutative diagrams we use Xy-pic
\usepackage[all]{xy}

% We use 2cell for 2-commutative diagrams.
\xyoption{2cell}
\UseAllTwocells

% We use multicol for the list of chapters between chapters
\usepackage{multicol}

% This is generally recommended for better output
\usepackage{lmodern}
\usepackage[T1]{fontenc}

% For cross-file-references

% Package for hypertext links:
\usepackage{hyperref}

% For any local file, say "hello.tex" you want to link to please

% Theorem environments.
%
\theoremstyle{plain}
\newtheorem{theorem}[subsection]{Theorem}
\newtheorem{proposition}[subsection]{Proposition}
\newtheorem{lemma}[subsection]{Lemma}

\theoremstyle{definition}
\newtheorem{definition}[subsection]{Definition}
\newtheorem{example}[subsection]{Example}
\newtheorem{exercise}[subsection]{Exercise}
\newtheorem{situation}[subsection]{Situation}

\theoremstyle{remark}
\newtheorem{remark}[subsection]{Remark}
\newtheorem{remarks}[subsection]{Remarks}

\numberwithin{equation}{subsection}

% Macros
%
\def\lim{\mathop{\mathrm{lim}}\nolimits}
\def\colim{\mathop{\mathrm{colim}}\nolimits}
\def\Spec{\mathop{\mathrm{Spec}}}
\def\Hom{\mathop{\mathrm{Hom}}\nolimits}
\def\Ext{\mathop{\mathrm{Ext}}\nolimits}
\def\SheafHom{\mathop{\mathcal{H}\!\mathit{om}}\nolimits}
\def\SheafExt{\mathop{\mathcal{E}\!\mathit{xt}}\nolimits}
\def\Sch{\mathit{Sch}}
\def\Mor{\mathop{\mathrm{Mor}}\nolimits}
\def\Ob{\mathop{\mathrm{Ob}}\nolimits}
\def\Sh{\mathop{\mathit{Sh}}\nolimits}
\def\NL{\mathop{N\!L}\nolimits}
\def\CH{\mathop{\mathrm{CH}}\nolimits}
\def\proetale{{pro\text{-}\acute{e}tale}}
\def\etale{{\acute{e}tale}}
\def\QCoh{\mathit{QCoh}}
\def\Ker{\mathop{\mathrm{Ker}}}
\def\Im{\mathop{\mathrm{Im}}}
\def\Coker{\mathop{\mathrm{Coker}}}
\def\Coim{\mathop{\mathrm{Coim}}}

% Boxtimes
%
\DeclareMathSymbol{\boxtimes}{\mathbin}{AMSa}{"02}

%
% Macros for moduli stacks/spaces
%
\def\QCohstack{\mathcal{QC}\!\mathit{oh}}
\def\Cohstack{\mathcal{C}\!\mathit{oh}}
\def\Spacesstack{\mathcal{S}\!\mathit{paces}}
\def\Quotfunctor{\mathrm{Quot}}
\def\Hilbfunctor{\mathrm{Hilb}}
\def\Curvesstack{\mathcal{C}\!\mathit{urves}}
\def\Polarizedstack{\mathcal{P}\!\mathit{olarized}}
\def\Complexesstack{\mathcal{C}\!\mathit{omplexes}}
% \Pic is the operator that assigns to X its picard group, usage \Pic(X)
% \Picardstack_{X/B} denotes the Picard stack of X over B
% \Picardfunctor_{X/B} denotes the Picard functor of X over B
\def\Pic{\mathop{\mathrm{Pic}}\nolimits}
\def\Picardstack{\mathcal{P}\!\mathit{ic}}
\def\Picardfunctor{\mathrm{Pic}}
\def\Deformationcategory{\mathcal{D}\!\mathit{ef}}

\hypersetup{hidelinks}
\widowpenalty=10000\clubpenalty=10000\displaywidowpenalty=10000
\begin{document}
\title[Stacks Project, Tag 0EJ9]{主理想形式层的代数化判据}
\maketitle
\noindent\textbf{Counting scope.} The primary problem is Tag 0EJ9, marked below. All preceding numbered statements are supporting author context and are not counted as additional problems. This is a local pilot with a temporary ID, not production admission.


\noindent Copyright (C) 2005--2025 Johan de Jong. Stacks Project authors. Source: \href{https://stacks.math.columbia.edu/tag/0EJ9}{Tag 0EJ9}.

\noindent Difficulty H2: 本题在有对偶化复形的Noetherian局部基底上，由(2,3)不等式证明主理想形式层可延拓至X。证明以局部上同调的有限模块控制障碍，构造有限分次A[T]模块并由齐次生成元推出Mittag–Leffler性质。代数化收口判据和局部对偶估计可引用，但障碍有限性与稳定像论证仍由本条完成，判H2。.

\noindent History: excerpt from revision local cached snapshot; exact chapter hashes; revision unknown, algebraization.tex, lines 6518--6654. Only reference presentation was adapted. Licensed under GNU FDL 1.2 or later; no invariant sections or cover texts. See accompanying COPYING.

\begin{situation}
\label{situation-algebraize}
Here $A$ is a Noetherian ring and $I \subset \mathfrak a \subset A$ are ideals.
We set $X = \Spec(A)$, $Y = V(I) = \Spec(A/I)$, and
$Z = V(\mathfrak a) = \Spec(A/\mathfrak a)$. Furthermore $U = X \setminus Z$.
\end{situation}

\noindent
Let $Y$ be a Noetherian scheme and let $Z \subset Y$ be a closed subset.
We define a function
\begin{equation}
\label{equation-delta-Z}
\delta^Y_Z = \delta_Z : Y \longrightarrow \mathbf{Z}_{\geq 0} \cup \{\infty\}
\end{equation}
which measures the ``distance'' of a point of $Y$ from $Z$.
For an informal discussion, please see Remark \href{https://stacks.math.columbia.edu/tag/0EJ0}{remark discussion (Tag 0EJ0)}.
Let $y \in Y$. We set $\delta_Z(y) = \infty$ if $y$ is contained
in a connected component of $Y$ which does not meet $Z$.
If $y$ is contained in a connected component of $Y$ which meets $Z$,
then we can find $k \geq 0$ and a system
$$
V_0 \subset W_0 \supset V_1 \subset W_1 \supset \ldots \supset
V_k \subset W_k
$$
of integral closed subschemes of $Y$ such that $V_0 \subset Z$
and $y \in W_k$ is the generic point. Set
$c_i = \text{codim}(V_i, W_i)$ for $i = 0, \ldots, k$
and $b_i = \text{codim}(V_{i + 1}, W_i)$ for $i = 0, \ldots, k - 1$.
For such a system we set
$$
\delta(V_0, W_0, V_1, \ldots, W_k) =
k +
\max_{i = 0, 1, \ldots, k}
(c_i + c_{i + 1} + \ldots + c_k - b_i - b_{i + 1} - \ldots - b_{k - 1})
$$
This is $\geq k$ as we can take $i = k$ and we have $c_k \geq 0$.
Finally, we set
$$
\delta_Z(y) = \min \delta(V_0, W_0, V_1, \ldots, W_k)
$$
where the minimum is over all systems of integral closed subschemes of $Y$
as above.

\begin{definition}
\label{definition-s-d-inequalities}
In Situation \ref{situation-algebraize} let $(\mathcal{F}_n)$ be an object
of $\textit{Coh}(U, I\mathcal{O}_U)$. Let $a, b$ be integers.
Let $\delta^Y_Z$ be as in (\ref{equation-delta-Z}).
We say
{\it $(\mathcal{F}_n)$ satisfies the $(a, b)$-inequalities} if for
$y \in U \cap Y$ and a prime $\mathfrak p \subset \mathcal{O}_{X, y}^\wedge$
with $\mathfrak p \not \in V(I\mathcal{O}_{X, y}^\wedge)$
\begin{enumerate}
\item if $V(\mathfrak p) \cap V(I\mathcal{O}_{X, y}^\wedge) \not =
\{\mathfrak m_y^\wedge\}$, then
$$
\text{depth}((\mathcal{F}^\wedge_y)_\mathfrak p) + \delta^Y_Z(y) \geq a
\quad\text{or}\quad
\text{depth}((\mathcal{F}^\wedge_y)_\mathfrak p) +
\dim(\mathcal{O}_{X, y}^\wedge/\mathfrak p) + \delta^Y_Z(y) > b
$$
\item if $V(\mathfrak p) \cap V(I\mathcal{O}_{X, y}^\wedge) =
\{\mathfrak m_y^\wedge\}$, then
$$
\text{depth}((\mathcal{F}^\wedge_y)_\mathfrak p) + \delta^Y_Z(y) > a
$$
\end{enumerate}
We say {\it $(\mathcal{F}_n)$ satisfies the strict $(a, b)$-inequalities}
if for $y \in U \cap Y$ and a prime
$\mathfrak p \subset \mathcal{O}_{X, y}^\wedge$ with
$\mathfrak p \not \in V(I\mathcal{O}_{X, y}^\wedge)$
we have
$$
\text{depth}((\mathcal{F}^\wedge_y)_\mathfrak p) + \delta^Y_Z(y) > a
\quad\text{or}\quad
\text{depth}((\mathcal{F}^\wedge_y)_\mathfrak p) +
\dim(\mathcal{O}_{X, y}^\wedge/\mathfrak p) + \delta^Y_Z(y) > b
$$
\end{definition}

\begin{lemma}
\label{lemma-helper-algebraize}
In Situation \ref{situation-algebraize} let $(M_n)$ be an inverse system of
$A$-modules as in Lemma \href{https://stacks.math.columbia.edu/tag/0DXT}{lemma system of modules (Tag 0DXT)} and let
$(\mathcal{F}_n)$ be the corresponding object of
$\textit{Coh}(U, I\mathcal{O}_U)$. Let $d \geq \text{cd}(A, I)$
and $s \geq 0$ be integers.
With notation as above assume
\begin{enumerate}
\item $A$ is local with maximal ideal $\mathfrak m = \mathfrak a$,
\item $A$ has a dualizing complex, and
\item $(\mathcal{F}_n)$ satisfies the $(s, s + d)$-inequalities
(Definition \ref{definition-s-d-inequalities}).
\end{enumerate}
Let $E$ be an injective hull of the residue field of $A$. Then for $i \leq s$
there exists a finite $A$-module $N$ annihilated by a power
of $I$ and for $n \gg 0$ compatible maps
$$
H^i_\mathfrak m(M_n) \to \Hom_A(N, E)
$$
whose cokernels are finite length $A$-modules and whose kernels $K_n$
form an inverse system such that $\Im(K_{n''} \to K_{n'})$ has finite
length for $n'' \gg n' \gg 0$.
\end{lemma}

\begin{proof}
Let $\omega_A^\bullet$ be a normalized dualizing complex. Then
$\delta^Y_Z = \delta$ is the dimension function associated with
this dualizing complex.
Observe that $\Ext^{-i}_A(M_n, \omega_A^\bullet)$ is a finite $A$-module
annihilated by $I^n$. Fix $0 \leq i \leq s$.
Below we will find $n_1 > n_0 > 0$ such that if we set
$$
N = \Im(\Ext^{-i}_A(M_{n_0}, \omega_A^\bullet) \to
\Ext^{-i}_A(M_{n_1}, \omega_A^\bullet))
$$
then the kernels of the maps
$$
N \to \Ext^{-i}_A(M_n, \omega_A^\bullet),\quad n \geq n_1
$$
are finite length $A$-modules and the cokernels $Q_n$ form a
system such that $\Im(Q_{n'} \to Q_{n''})$ has finite length
for $n'' \gg n' \gg n_1$. This is equivalent to the statement that
the system $\{\Ext^{-i}_A(M_n, \omega_A^\bullet)\}_{n \geq 1}$
is essentially constant in the quotient of the category of finite
$A$-modules modulo the Serre subcategory of finite length $A$-modules.
By the local duality theorem
(Dualizing Complexes, Lemma \href{https://stacks.math.columbia.edu/tag/0AAK}{lemma special case local duality (Tag 0AAK)})
and Matlis duality
(Dualizing Complexes, Proposition \href{https://stacks.math.columbia.edu/tag/08Z9}{proposition matlis (Tag 08Z9)})
we conclude that there are maps
$$
H^i_\mathfrak m(M_n) \to \Hom_A(N, E),\quad n \geq n_1
$$
as in the statement of the lemma.

\medskip\noindent
Pick $f \in \mathfrak m$. Let $B = A_f^\wedge$ be the $I$-adic completion
of the localization $A_f$. Recall that
$\omega_{A_f}^\bullet = \omega_A^\bullet \otimes_A A_f$
and $\omega_B^\bullet = \omega_A^\bullet \otimes_A B$ are dualizing
complexes (Dualizing Complexes, Lemma \href{https://stacks.math.columbia.edu/tag/0A7G}{lemma dualizing localize (Tag 0A7G)}
and \href{https://stacks.math.columbia.edu/tag/0DWD}{lemma completion henselization dualizing (Tag 0DWD)}).
Let $M$ be the finite $B$-module $\lim M_{n, f}$ (compare with
discussion in Cohomology of Schemes, Lemma
\href{https://stacks.math.columbia.edu/tag/087W}{lemma inverse systems affine (Tag 087W)}). Then
$$
\Ext^{-i}_A(M_n, \omega_A^\bullet)_f =
\Ext^{-i}_{A_f}(M_{n, f}, \omega_{A_f}^\bullet) =
\Ext^{-i}_B(M/I^n M, \omega_B^\bullet)
$$
Since $\mathfrak m$ can be generated by finitely many $f \in \mathfrak m$
it suffices to show that for each $f$ the system
$$
\{\Ext^{-i}_B(M/I^n M, \omega_B^\bullet)\}_{n \geq 1}
$$
is essentially constant. Some details omitted.

\medskip\noindent
Let $\mathfrak q \subset IB$ be a prime ideal. Then $\mathfrak q$ corresponds
to a point $y \in U \cap Y$. Observe that
$\delta(\mathfrak q) = \dim(\overline{\{y\}})$
is also the value of the dimension function associated to $\omega_B^\bullet$
(we omit the details; use that $\omega_B^\bullet$ is gotten from
$\omega_A^\bullet$ by tensoring up with $B$). Assumption
(3) guarantees via Lemma \href{https://stacks.math.columbia.edu/tag/0EJ4}{lemma elementary (Tag 0EJ4)}
that Lemma \href{https://stacks.math.columbia.edu/tag/0EFR}{lemma algebraize local cohomology bis (Tag 0EFR)}
applies to
$B_\mathfrak q, IB_\mathfrak q, \mathfrak qB_\mathfrak q, M_\mathfrak q$
with $s$ replaced by $s - \delta(y)$. We obtain that
$$
H^{i - \delta(\mathfrak q)}_{\mathfrak qB_\mathfrak q}(M_\mathfrak q) =
\lim H^{i - \delta(\mathfrak q)}_{\mathfrak qB_\mathfrak q}(
(M/I^nM)_\mathfrak q)
$$
and this module is annihilated by a power of $I$.
By Lemma \href{https://stacks.math.columbia.edu/tag/0EHB}{lemma terrific (Tag 0EHB)} we find that the inverse systems
$H^{i - \delta(\mathfrak q)}_{\mathfrak qB_\mathfrak q}((M/I^nM)_\mathfrak q)$
are essentially constant with value
$H^{i - \delta(\mathfrak q)}_{\mathfrak qB_\mathfrak q}(M_\mathfrak q)$.
Since $(\omega_B^\bullet)_\mathfrak q[-\delta(\mathfrak q)]$ is a normalized
dualizing complex on $B_\mathfrak q$ the local duality theorem
shows that the system
$$
\Ext^{-i}_B(M/I^n M, \omega_B^\bullet)_\mathfrak q
$$
is essentially constant with value
$\Ext^{-i}_B(M, \omega_B^\bullet)_\mathfrak q$.

\medskip\noindent
To finish the proof we globalize as in the proof of
Lemma \href{https://stacks.math.columbia.edu/tag/0EFX}{lemma bootstrap (Tag 0EFX)}; the argument here is easier
because we know the value of our system already. Namely, consider the maps
$$
\alpha_n :
\Ext^{-i}_B(M/I^n M, \omega_B^\bullet)
\longrightarrow
\Ext^{-i}_B(M, \omega_B^\bullet)
$$
for varying $n$. By the above, for every $\mathfrak q$ we can find an
$n$ such that $\alpha_n$ is surjective after localization at $\mathfrak q$.
Since $B$ is Noetherian and $\Ext^{-i}_B(M, \omega_B^\bullet)$
a finite module, we can find an $n$ such that $\alpha_n$ is surjective.
For any $n$ such that $\alpha_n$ is surjective, given a prime
$\mathfrak q \in V(IB)$ we can find an $n' > n$ such that
$\Ker(\alpha_n)$ maps to zero in $\Ext^{-i}(M/I^{n'}M, \omega_B^\bullet)$
at least after localizing at $\mathfrak q$.
Since $\Ker(\alpha_n)$ is a finite $A$-module and since supports of
sections are quasi-compact, we can find an $n'$ such that
$\Ker(\alpha_n)$ maps to zero in $\Ext^{-i}(M/I^{n'}M, \omega_B^\bullet)$.
In this way we see that $\Ext^{-i}(M/I^n M, \omega_B^\bullet)$
is essentially constant with value $\Ext^{-i}(M, \omega_B^\bullet)$.
This finishes the proof.
\end{proof}


\begin{lemma}
\label{lemma-when-done}
In Situation \ref{situation-algebraize} let $(\mathcal{F}_n)$ be an object
of $\textit{Coh}(U, I\mathcal{O}_U)$. Let $M$ be as in (\href{https://stacks.math.columbia.edu/tag/0EHD}{equation guess (Tag 0EHD)}).
Assume
\begin{enumerate}
\item[(a)] the inverse system $H^0(U, \mathcal{F}_n)$ has Mittag-Leffler,
\item[(b)] the limit topology on $M$ agrees with the $I$-adic topology, and
\item[(c)] the image of $M \to H^0(U, \mathcal{F}_n)$ is a finite $A$-module
for all $n$.
\end{enumerate}
Then $(\mathcal{F}_n)$ extends canonically to $X$.
In particular, if $A$ is $I$-adically complete, then
$(\mathcal{F}_n)$ is the completion of a coherent $\mathcal{O}_U$-module.
\end{lemma}

\begin{proof}
Since $H^0(U, \mathcal{F}_n)$ has the Mittag-Leffler condition
and since the limit topology on $M$ is the $I$-adic topology
we see that $\{M/I^nM\}$ and $\{H^0(U, \mathcal{F}_n)\}$
are pro-isomorphic inverse systems of $A$-modules.
Thus if we set
$$
\mathcal{G}_n = \widetilde{M/I^n M}
$$
then we see that to verify the condition in
Definition \ref{definition-canonically-algebraizable}
it suffices to show that $M$ is a finite module over the
$I$-adic completion of $A$. This follows from
the fact that $M/I^n M$ is finite by condition (c)
and the above and
Algebra, Lemma \href{https://stacks.math.columbia.edu/tag/031D}{lemma finite over complete ring (Tag 031D)}.
\end{proof}


\begin{definition}
\label{definition-canonically-algebraizable}
In Situation \ref{situation-algebraize} let $(\mathcal{F}_n)$ be an
object of $\textit{Coh}(U, I\mathcal{O}_U)$. We say
{\it $(\mathcal{F}_n)$ canonically extends to $X$} if the the
inverse system
$$
\{\widetilde{H^0(U, \mathcal{F}_n)}\}_{n \geq 1}
$$
in $\QCoh(\mathcal{O}_X)$ is pro-isomorphic to an object
$(\mathcal{G}_n)$ of $\textit{Coh}(X, I\mathcal{O}_X)$.
\end{definition}

\section*{Primary problem: Tag 0EJ9}
\begin{lemma}
\label{lemma-algebraization-principal-bis}
In Situation \ref{situation-algebraize} let $(\mathcal{F}_n)$ be an object
of $\textit{Coh}(U, I\mathcal{O}_U)$. Assume
\begin{enumerate}
\item $A$ is local and $\mathfrak a = \mathfrak m$ is the maximal ideal,
\item $A$ has a dualizing complex,
\item $I = (f)$ is a principal ideal,
\item $(\mathcal{F}_n)$ satisfies the $(2, 3)$-inequalities.
\end{enumerate}
Then $(\mathcal{F}_n)$ extends to $X$. In particular, if $A$ is
$I$-adically complete, then $(\mathcal{F}_n)$ is the completion
of a coherent $\mathcal{O}_U$-module. 
\end{lemma}

\begin{proof}
Recall that $\textit{Coh}(U, I\mathcal{O}_U)$ is an abelian category, see
Cohomology of Schemes, Lemma \href{https://stacks.math.columbia.edu/tag/087X}{lemma inverse systems abelian (Tag 087X)}.
Over affine opens of $U$ the object $(\mathcal{F}_n)$
corresponds to a finite module over a Noetherian ring
(Cohomology of Schemes, Lemma \href{https://stacks.math.columbia.edu/tag/087W}{lemma inverse systems affine (Tag 087W)}).
Thus the kernels of the maps $f^N : (\mathcal{F}_n) \to (\mathcal{F}_n)$
stabilize for $N$ large enough. By
Lemmas \href{https://stacks.math.columbia.edu/tag/0EIU}{lemma map kernel cokernel on closed (Tag 0EIU)} and
\href{https://stacks.math.columbia.edu/tag/0EIL}{lemma essential image completion (Tag 0EIL)}
in order to prove the lemma
we may replace $(\mathcal{F}_n)$ by the image of such a map.
Thus we may assume $f$ is injective on $(\mathcal{F}_n)$.
After this replacement the equivalent conditions of
Lemma \href{https://stacks.math.columbia.edu/tag/0EH9}{lemma equivalent f good (Tag 0EH9)} hold for the inverse system
$(\mathcal{F}_n)$ on $U$. We will use this without further mention
in the rest of the proof.

\medskip\noindent
We will check hypotheses (a), (b), and (c) of
Lemma \ref{lemma-when-done}. Hypothesis (b) holds by
Cohomology, Lemma \href{https://stacks.math.columbia.edu/tag/0EHA}{lemma topology I adic f (Tag 0EHA)}.

\medskip\noindent
Pick an inverse system of modules $\{M_n\}$ as in
Lemma \href{https://stacks.math.columbia.edu/tag/0DXT}{lemma system of modules (Tag 0DXT)}.
We may assume $H^0_\mathfrak m(M_n) = 0$ by replacing $M_n$ by
$M_n/H^0_\mathfrak m(M_n)$ if necessary. Then we obtain short exact
sequences
$$
0 \to M_n \to H^0(U, \mathcal{F}_n) \to H^1_\mathfrak m(M_n) \to 0
$$
for all $n$. Let $E$ be an injective hull of the residue field of $A$.
By Lemma \ref{lemma-helper-algebraize} and our current assumption (4)
we can choose, an integer $m \geq 0$, finite $A$-modules
$N_1$ and $N_2$ annihilated by $f^c$ for some $c \geq 0$ and
compatible systems of maps
$$
H^i_\mathfrak m(M_n) \to \Hom_A(N_i, E), \quad i = 1, 2
$$
for $n \geq m$
with the properties stated in the lemma.

\medskip\noindent
We know that $M = \lim H^0(U, \mathcal{F}_n)$ is an $A$-module whose
limit topology is the $f$-adic topology. Thus, given $n$, the module
$M/f^nM$ is a subquotient of $H^0(U, \mathcal{F}_N)$ for some $N \gg n$.
Looking at the information obtained above we see that
$f^cM/f^nM$ is a finite $A$-module. Since $f$ is a nonzerodivisor
on $M$ we conclude that $M/f^{n - c}M$ is a finite $A$-module.
In this way we see that hypothesis (c) of Lemma \ref{lemma-when-done} holds.

\medskip\noindent
Next, we study the module
$$
Ob = \lim H^1(U, \mathcal{F}_n) = \lim H^2_\mathfrak m(M_n)
$$
For $n \geq m$ let $K_n$ be the kernel of the map
$H^2_\mathfrak m(M_n) \to \Hom_A(N_2, E)$.
Set $K = \lim K_n$. We obtain an exact sequence
$$
0 \to K \to Ob \to \Hom_A(N_2, E)
$$
By the above the limit topology on $Ob = \lim H^2_\mathfrak m(M_n)$
is the $f$-adic topology. Since $N_2$ is annihilated by $f^c$
we conclude the same is true for the limit topology on $K = \lim K_n$.
Thus $K/fK$ is a subquotient of $K_n$ for $n \gg 1$.
However, since $\{K_n\}$ is pro-isomorphic to an inverse system of
finite length $A$-modules (by the conclusion of
Lemma \ref{lemma-helper-algebraize})
we conclude that $K/fK$ is a subquotient of a finite length
$A$-module. It follows that $K$ is a finite $A$-module, see
Algebra, Lemma \href{https://stacks.math.columbia.edu/tag/031D}{lemma finite over complete ring (Tag 031D)}.
(In fact, we even see that $\dim(\text{Supp}(K)) = 1$ but
we will not need this.)

\medskip\noindent
Given $n \geq 1$ consider the boundary map
$$
\delta_n :
H^0(U, \mathcal{F}_n)
\longrightarrow
\lim_N H^1(U, f^n\mathcal{F}_N) \xrightarrow{f^{-n}} Ob
$$
(the second map is an isomorphism)
coming from the short exact sequences
$$
0 \to f^n\mathcal{F}_N \to \mathcal{F}_N \to \mathcal{F}_n \to 0
$$
For each $n$ set
$$
P_n = \Im(H^0(U, \mathcal{F}_{n + m}) \to H^0(U, \mathcal{F}_n))
$$
where $m$ is as above. Observe that $\{P_n\}$ is an inverse
system and that the map $f : \mathcal{F}_n \to \mathcal{F}_{n + 1}$
on global sections maps $P_n$ into $P_{n + 1}$.
If $p \in P_n$, then $\delta_n(p) \in K \subset Ob$
because $\delta_n(p)$ maps to zero in
$H^1(U, f^n\mathcal{F}_{n + m}) = H^2_\mathfrak m(M_m)$
and the composition of $\delta_n$ and $Ob \to \Hom_A(N_2, E)$
factors through $H^2_\mathfrak m(M_m)$ by our choice of $m$.
Hence
$$
\bigoplus\nolimits_{n \geq 0} \Im(P_n \to Ob)
$$
is a finite graded $A[T]$-module where $T$ acts via multiplication by $f$.
Namely, it is a graded submodule of $K[T]$ and $K$ is finite over $A$.
Arguing as in the proof of
Cohomology, Lemma \href{https://stacks.math.columbia.edu/tag/0GYK}{lemma ML general (Tag 0GYK)}\footnote{Choose
homogeneous generators of the form $\delta_{n_j}(p_j)$ for the displayed
module. Then if $k = \max(n_j)$ we find that for $n \geq k$
and any $p \in P_n$ we can find $a_j \in A$ such that
$p - \sum a_j f^{n - n_j} p_j$ is in the kernel of $\delta_n$
and hence in the image of $P_{n'}$ for all $n' \geq n$.
Thus $\Im(P_n \to P_{n - k}) = \Im(P_{n'} \to P_{n - k})$
for all $n' \geq n$.}
we find that the inverse system $\{P_n\}$ satisfies ML.
Since $\{P_n\}$ is pro-isomorphic to $\{H^0(U, \mathcal{F}_n)\}$
we conclude that $\{H^0(U, \mathcal{F}_n)\}$ has ML.
Thus hypothesis (a) of Lemma \ref{lemma-when-done}
holds and the proof is complete.
\end{proof}
\end{document}
